% ======================================================================
% capitulos/capitulo3.tex
% ----------------------------------------------------------------------
% ENTREGABLE 3 — PLAN DE GESTIÓN DEL PROYECTO
% Empresa caso de estudio: WISP Ingeniería en Redes y Telecomunicaciones
% Enfoque híbrido: PMBOK (PMI) + Scrum
% Contenido trasladado desde Entregable3_PlanGestion_GTI_WISP.docx
% ======================================================================
\chapter{Plan de Gestión del Proyecto}
\label{chap:plangestion}

% ----------------------------------------------------------------------
% Descripción del producto (introducción del capítulo)
% ----------------------------------------------------------------------
El presente documento constituye el Plan de Gestión del Proyecto para el desarrollo
del \textbf{Sistema de Gestión Integral para ISP de Fibra Óptica Multi-Zona} de la
empresa WISP. Es el instrumento técnico que describe cómo será planificado,
ejecutado, monitoreado y controlado el proyecto, integrando la definición del
alcance, el cronograma, el presupuesto, la gestión de riesgos, la gestión de
interesados y los mecanismos de control.

Este plan adopta un enfoque \textbf{híbrido}, reconocido por la Guía PMBOK 7.ª
edición (PMI). La planificación general ---acta, alcance, cronograma, presupuesto y
línea base--- sigue el enfoque predictivo del PMI, mientras que el desarrollo del
software se ejecuta bajo un enfoque adaptativo (Scrum: \textit{backlog}, sprints e
incrementos). Esta combinación asegura disciplina metodológica y trazabilidad de
resultados, adecuada a la naturaleza mixta del proyecto (documentación de gestión y
construcción iterativa de software).

\begin{upeubox}{Alcance temporal del proyecto integrador}
El proyecto integrador abarca las tres líneas (GTI, Software e Infraestructura) y se
organiza en dos semestres:
\begin{itemize}[leftmargin=1.4em]
  \item \textbf{Semestre 1 (16/03/2026 -- 30/06/2026):} fase de iniciación y
  planificación. Se completa íntegramente la línea de GTI y se produce la
  documentación base de Infraestructura y Software (diseño, SRS, arquitectura).
  \item \textbf{Semestre 2 (agosto -- inicios de diciembre 2026):} fase de ejecución.
  Se implementan y despliegan las soluciones de Infraestructura y Software. La GTI,
  ya cerrada, gobierna el proyecto integrador.
\end{itemize}
\end{upeubox}

% ======================================================================
\section{Acta de Constitución del Proyecto}
\label{sec:acta}

El Acta de Constitución (\textit{Project Charter}) es el documento que autoriza
formalmente el proyecto, nombra al director del proyecto y establece los objetivos,
el alcance preliminar, los interesados clave, las restricciones y los supuestos. Es
la salida del proceso «Desarrollar el Acta de Constitución» del Grupo de Procesos de
Inicio (PMBOK).

\begin{center}
\captionof{table}{Acta de constitución del proyecto.}
\label{tab:acta}
\begin{tabularx}{\textwidth}{L{4.2cm} X}
\toprule
\textbf{Campo} & \textbf{Contenido} \\
\midrule
Nombre del proyecto & Sistema de Gestión Integral para ISP de Fibra Óptica
Multi-Zona -- WISP. \\
Sponsor (patrocinador) & Gerencia de WISP Ingeniería en Redes y Telecomunicaciones
(asume la infraestructura de producción y valida los entregables). \\
Product Owner & Cristian Cabana Sulca (conoce el negocio; define y prioriza los
requerimientos). \\
Scrum Master & Jherson Jean Piero Fernández Sánchez (facilita el proceso, remueve
impedimentos). \\
Equipo de desarrollo & Alessandro Pastor Mamani Mamani (backend y móvil) y Andrés
Lino Montes Mamani (backend y frontend). \\
Fecha de inicio & 16 de marzo de 2026. \\
Fecha de cierre (GTI) & 30 de junio de 2026 (fin del semestre 1). \\
Justificación & WISP opera con procesos manuales y un sistema heredado deficiente, lo
que limita su eficiencia y escalabilidad (según diagnóstico y Business Case previos). \\
\bottomrule
\end{tabularx}
\end{center}

\subsection{Objetivos del proyecto}

\begin{center}
\captionof{table}{Objetivos del proyecto y métricas de éxito.}
\label{tab:objetivos-acta}
\begin{tabularx}{\textwidth}{L{2.2cm} X L{4cm}}
\toprule
\textbf{Tipo} & \textbf{Objetivo} & \textbf{Métrica de éxito} \\
\midrule
Alcance & Entregar el sistema web y app móvil integrado con los 13 MikroTik. &
100\% de módulos y entregables definidos. \\
Tiempo & Completar la documentación de gestión (GTI) en el semestre 1. &
Cierre al 30/06/2026. \\
Costo & Mantener la inversión de WISP dentro del presupuesto aprobado. &
$\leq$ S/ 1\,000 de inversión inicial. \\
Calidad & Cumplir criterios de calidad ISO/IEC~25010 en el software. &
Aprobación de pruebas y jurado. \\
\bottomrule
\end{tabularx}
\end{center}

\subsection{Alcance preliminar}

\begin{itemize}[leftmargin=1.4em]
  \item \textbf{Incluye:} sistema web de administración (Django + Angular +
  PostgreSQL), app móvil con roles de técnico y vendedor, integración con los 13
  MikroTik (WireGuard + API RouterOS), migración de datos y automatización de
  notificaciones WhatsApp.
  \item \textbf{No incluye:} automatización del sistema de pagos (Yape/Plin),
  provisión automática de ONU en la OLT, ni el módulo contable/facturación
  electrónica (fuera del alcance de esta fase).
\end{itemize}

\subsection{Restricciones y supuestos}

\begin{center}
\captionof{table}{Restricciones y supuestos del proyecto.}
\label{tab:restricciones}
\begin{tabularx}{\textwidth}{X X}
\toprule
\textbf{Restricciones} & \textbf{Supuestos} \\
\midrule
El proyecto documental de GTI debe cerrar en el semestre 1 (30/06/2026). &
WISP proveerá el VPS, dominio y accesos a los MikroTik a tiempo. \\
El equipo desarrolla en paralelo a su carga académica. &
Los 13 MikroTik usan el mismo modelo (RB4011) y RouterOS~7 (homogéneo). \\
Inversión de WISP limitada (MYPE). &
El personal de WISP colaborará en la validación y capacitación. \\
Conectividad variable en zonas rurales. &
La infraestructura de red actual se mantiene operativa durante el proyecto. \\
\bottomrule
\end{tabularx}
\end{center}

% ======================================================================
\section{Gestión del Alcance}
\label{sec:alcance}

La Estructura de Desglose del Trabajo (EDT/WBS) es una descomposición jerárquica del
trabajo total del proyecto en entregables y paquetes de trabajo. Cumple la
\textbf{regla del 100\,\%}: la suma de los niveles inferiores representa la totalidad
del trabajo del nivel superior, sin trabajo adicional ni faltante.

\subsection{Estructura de Desglose del Trabajo (EDT / WBS)}

La EDT del proyecto se organiza en cinco entregables de primer nivel, cada uno
descompuesto en paquetes de trabajo:

\begin{center}
\captionof{table}{Estructura de Desglose del Trabajo (EDT/WBS).}
\label{tab:edt}
\begin{tabularx}{\textwidth}{L{1.6cm} X}
\toprule
\textbf{Código} & \textbf{Entregable / Paquete de trabajo} \\
\midrule
\textbf{1} & \textbf{GESTIÓN DEL PROYECTO (GTI)} \\
1.1 & Diagnóstico organizacional y alineamiento estratégico. \\
1.2 & Business Case (caso de negocio). \\
1.3 & Plan de gestión del proyecto. \\
1.4 & Mejora de procesos (AS-IS / TO-BE). \\
1.5 & Solución TIC integrada (arquitectura). \\
\addlinespace
\textbf{2} & \textbf{DEFINICIÓN DEL SISTEMA (Software E1)} \\
2.1 & SRS (IEEE~29148): requerimientos funcionales y no funcionales. \\
2.2 & Prototipos navegables (Figma) validados. \\
2.3 & Documento de arquitectura (C4 / IEEE~42010) y modelado UML. \\
\addlinespace
\textbf{3} & \textbf{BASE DE DATOS (Software E2)} \\
3.1 & Modelo de datos (ER + lógico + diccionario). \\
3.2 & Scripts SQL (DDL + DML), procedimientos, seguridad. \\
\addlinespace
\textbf{4} & \textbf{CONSTRUCCIÓN DEL SISTEMA (Software E3 -- Semestre 2)} \\
4.1 & Backend Django (API REST) e integración MikroTik (WireGuard + RouterOS API). \\
4.2 & Frontend web Angular (administración). \\
4.3 & App móvil (roles técnico y vendedor). \\
4.4 & Módulo WhatsApp (notificaciones) y despliegue en nube. \\
\addlinespace
\textbf{5} & \textbf{VALIDACIÓN Y CIERRE (Software E4-E5 -- Semestre 2)} \\
5.1 & Pruebas, CI/CD, auditoría de calidad. \\
5.2 & Video pitch y sustentación ante jurado. \\
\bottomrule
\end{tabularx}
\end{center}

\begin{upeubox}{Foco de este plan}
Este Plan de Gestión pertenece a la línea GTI, que se ejecuta y cierra en el semestre
1. Por ello, el cronograma, el presupuesto y los sprints detallados a continuación se
centran en el paquete~1 (Gestión del Proyecto) y en la documentación base de los
paquetes~2 y~3, que son las salidas del semestre~1. Los paquetes~4 y~5 (construcción
y validación) corresponden al semestre~2 y se planifican a nivel de hitos.
\end{upeubox}

\subsection{Diccionario de la EDT (entregables principales)}

El diccionario de la EDT describe cada entregable de primer nivel, su descripción y
su criterio de aceptación:

\begin{center}
\captionof{table}{Diccionario de la EDT.}
\label{tab:diccionario-edt}
\begin{tabularx}{\textwidth}{L{2.9cm} X X}
\toprule
\textbf{Entregable} & \textbf{Descripción} & \textbf{Criterio de aceptación} \\
\midrule
1. Gestión (GTI) & Documentación completa de diagnóstico, caso de negocio, plan,
procesos y arquitectura. & Aprobación de los 4 componentes por el jurado. \\
2. Definición del sistema & SRS, prototipos, arquitectura y UML del sistema. &
SRS validado y prototipos aprobados por WISP. \\
3. Base de datos & Modelo, scripts y seguridad de la base de datos. &
BD normalizada y scripts ejecutables. \\
4. Construcción & Sistema web y móvil funcional e integrado con MikroTik. &
Sistema operativo y desplegado (semestre 2). \\
5. Validación y cierre & Pruebas, CI/CD y sustentación final. &
Aprobación del jurado (semestre 2). \\
\bottomrule
\end{tabularx}
\end{center}

% ======================================================================
\section{Gestión del Cronograma}
\label{sec:cronograma}

El cronograma se construye definiendo las actividades, secuenciándolas, estimando su
duración y desarrollando el modelo de programación. Se representa mediante un diagrama
de Gantt, con hitos que marcan eventos significativos del proyecto (PMBOK, Área de
Conocimiento: Gestión del Cronograma).

\subsection{Lista de actividades}

Las actividades del semestre~1 (línea GTI), organizadas por fase:

\begin{center}
\captionof{table}{Lista de actividades del semestre 1.}
\label{tab:actividades}
\begin{tabularx}{\textwidth}{C{1cm} X L{2.4cm} L{2.6cm}}
\toprule
\textbf{ID} & \textbf{Actividad} & \textbf{Duración} & \textbf{Responsable} \\
\midrule
A1 & Inducción y conformación de equipos & Semana 1-2 & Equipo \\
A2 & Diagnóstico organizacional y alineamiento estratégico & 2 semanas & Product Owner \\
A3 & Business Case (caso de negocio) & 2 semanas & Product Owner \\
A4 & Plan de gestión del proyecto & 2 semanas & Scrum Master \\
A5 & Modelado de procesos AS-IS / TO-BE & 2 semanas & PO + equipo \\
A6 & SRS y prototipos del sistema & 3 semanas & Developers \\
A7 & Arquitectura de solución TIC + UML & 2 semanas & Developers \\
A8 & Diseño de base de datos & 2 semanas & Developers \\
A9 & Consolidación y sustentación (semestre 1) & 2 semanas & Equipo \\
\bottomrule
\end{tabularx}
\end{center}

\subsection{Diagrama de Gantt (Semestre 1: 16/03 -- 30/06/2026)}

Cada columna representa una semana. Las celdas coloreadas indican el período de
ejecución de cada actividad:

\begin{center}
\captionof{table}{Diagrama de Gantt del semestre 1 (15 semanas).}
\label{tab:gantt}
\setlength{\tabcolsep}{3pt}
\renewcommand{\arraystretch}{1.15}
{\footnotesize
\begin{tabular}{l *{15}{c}}
\toprule
\textbf{Act.} & \textbf{S1} & \textbf{S2} & \textbf{S3} & \textbf{S4} & \textbf{S5}
& \textbf{S6} & \textbf{S7} & \textbf{S8} & \textbf{S9} & \textbf{S10} & \textbf{S11}
& \textbf{S12} & \textbf{S13} & \textbf{S14} & \textbf{S15} \\
\midrule
A1 Inducción      & \gc & \gc &     &     &     &     &     &     &     &     &     &     &     &     &     \\
A2 Diagnóstico    &     & \gc & \gc &     &     &     &     &     &     &     &     &     &     &     &     \\
A3 Business Case  &     &     & \gc & \gc &     &     &     &     &     &     &     &     &     &     &     \\
A4 Plan gestión   &     &     &     & \gc & \gc & \gc &     &     &     &     &     &     &     &     &     \\
A5 Procesos       &     &     &     &     &     & \gc & \gc &     &     &     &     &     &     &     &     \\
A6 SRS+Prototipos &     &     &     &     &     &     &     & \gc & \gc & \gc &     &     &     &     &     \\
A7 Arquitect.+UML &     &     &     &     &     &     &     &     &     & \gc & \gc & \gc &     &     &     \\
A8 Base datos     &     &     &     &     &     &     &     &     &     &     &     & \gc & \gc &     &     \\
A9 Consol.+Sust.  &     &     &     &     &     &     &     &     &     &     &     &     &     & \gc & \gc \\
\bottomrule
\end{tabular}}
\end{center}

\subsection{Hitos principales}

\begin{center}
\captionof{table}{Hitos principales del semestre 1.}
\label{tab:hitos}
\begin{tabularx}{\textwidth}{C{1cm} X L{4cm}}
\toprule
\textbf{Hito} & \textbf{Descripción} & \textbf{Fecha estimada} \\
\midrule
H1 & Acta de conformidad e inducción aprobada & Semana 2 (fin marzo) \\
H2 & Revisión de avance 1 (diagnóstico + business case) & Semana 5-6 (abril) \\
H3 & Brief aprobado (problema y alcance validados) & Semana 8 (mayo) \\
H4 & Revisión de avance 2 (procesos + arquitectura) & Semana 11-12 (junio) \\
H5 & Sustentación final semestre 1 (GTI cerrada) & Semana 15-16 (30/06) \\
\bottomrule
\end{tabularx}
\end{center}

% ======================================================================
\section{Gestión de Costos}
\label{sec:costos-plan}

El presupuesto es la suma de los costos estimados de las actividades, agregada para
establecer la \textbf{línea base de costos}, contra la cual se mide y controla el
desempeño del proyecto (PMBOK, Gestión de los Costos).

\subsection{Presupuesto detallado}

El presupuesto integra la inversión que asume WISP (única) y los costos operativos del
primer año. Se retoma la estimación del Business Case (Entregable~2), organizada aquí
por categorías de gestión de costos:

\begin{center}
\captionof{table}{Presupuesto detallado del proyecto (soles).}
\label{tab:presupuesto}
\begin{tabularx}{\textwidth}{L{3.2cm} X R{2.2cm}}
\toprule
\textbf{Categoría} & \textbf{Concepto} & \textbf{Monto (S/)} \\
\midrule
Infraestructura & VPS en la nube (setup + 1.er mes) & 150 \\
Infraestructura & Dominio web + SSL + correo corporativo & 150 \\
Implementación & Migración de datos y puesta en marcha & 300 \\
Gestión del cambio & Capacitación al personal (2 sesiones) & 200 \\
Reserva & Contingencia / imprevistos ($\sim$10\%) & 200 \\
\midrule
\textbf{INVERSIÓN INICIAL} & Desembolso único de WISP al arranque & \textbf{1\,000} \\
Operación (anual) & VPS mensual + dominio + mantenimiento + backup & 3\,780 \\
\textbf{TOTAL AÑO 1} & Inversión inicial + operación del primer año & \textbf{4\,780} \\
\bottomrule
\end{tabularx}
\end{center}

\begin{upeubox}{Costo del desarrollo (recurso humano del equipo)}
El esfuerzo de desarrollo del equipo de 4 integrantes no representa un costo monetario
para WISP (es el aporte del proyecto de fin de carrera), pero su valor de mercado se
estima en \textbf{S/ 30\,000}. Se registra como costo asumido por el equipo para
efectos de trazabilidad y cálculo del retorno, conforme se detalló en el Business Case.
\end{upeubox}

\subsection{Línea base de costos}

La línea base de costos distribuye el gasto acumulado a lo largo del tiempo (curva~S).
Para el semestre~1, la inversión se concentra hacia el final (puesta en marcha),
mientras que los costos operativos inician tras el despliegue:

\begin{center}
\captionof{table}{Línea base de costos (curva S).}
\label{tab:linea-base}
\begin{tabularx}{\textwidth}{X R{3.2cm} R{3.2cm}}
\toprule
\textbf{Momento} & \textbf{Costo del período (S/)} & \textbf{Costo acumulado (S/)} \\
\midrule
Inicio (planificación) & 0 & 0 \\
Setup infraestructura & 300 & 300 \\
Puesta en marcha + capacitación & 700 & 1\,000 \\
Operación (resto del año) & 3\,780 & 4\,780 \\
\bottomrule
\end{tabularx}
\end{center}

% ======================================================================
\section{Gestión de Riesgos}
\label{sec:riesgos-plan}

La gestión de riesgos comprende identificar los riesgos, realizar el análisis
cualitativo (evaluando probabilidad e impacto para priorizarlos) y planificar la
respuesta a los riesgos (estrategias: evitar, mitigar, transferir o aceptar), conforme
al PMBOK.

\subsection{Identificación y análisis cualitativo}

Se retoman y amplían los riesgos del Business Case, evaluándolos por probabilidad e
impacto (escala 1-5) y priorizándolos por severidad ($P\times I$):

\begin{center}
\captionof{table}{Identificación y análisis cualitativo de riesgos.}
\label{tab:riesgos-identificacion}
\begin{tabularx}{\textwidth}{C{1cm} X C{1.1cm} C{1.1cm} C{1.1cm} C{1.6cm}}
\toprule
\textbf{ID} & \textbf{Riesgo} & \textbf{Prob} & \textbf{Imp} & \textbf{Sev} &
\textbf{Nivel} \\
\midrule
R1 & Baja adopción del sistema por el personal. & 4 & 4 & 16 & Alto \\
R2 & Fallo en la integración con MikroTik vía API. & 3 & 4 & 12 & Medio \\
R3 & Pérdida de datos en la migración. & 3 & 5 & 15 & Alto \\
R4 & Conectividad inestable en zonas rurales. & 4 & 3 & 12 & Medio \\
R5 & Retrasos por carga académica del equipo. & 4 & 3 & 12 & Medio \\
R6 & Caída del VPS en la nube. & 2 & 4 & 8 & Medio \\
R7 & Brechas de seguridad / acceso no autorizado. & 3 & 5 & 15 & Alto \\
\bottomrule
\end{tabularx}
\end{center}

\subsection{Plan de respuesta a los riesgos}

\begin{center}
\captionof{table}{Plan de respuesta a los riesgos.}
\label{tab:riesgos-respuesta}
\begin{tabularx}{\textwidth}{C{1cm} L{2.4cm} X}
\toprule
\textbf{ID} & \textbf{Estrategia} & \textbf{Plan de respuesta} \\
\midrule
R1 & Mitigar & Capacitación práctica, interfaz simple tipo app conocida,
acompañamiento inicial y respaldo del administrador. \\
R2 & Mitigar & Pruebas de integración por zona; los 13 MikroTik son homogéneos
(RB4011, RouterOS~7); ambiente de prueba previo. \\
R3 & Evitar / Mitigar & Respaldo previo, migración por etapas, validación de datos y
conservación del sistema heredado hasta confirmar consistencia. \\
R4 & Mitigar & Apps con modo offline / sincronización diferida para operar en campo
sin señal. \\
R5 & Mitigar & Metodología ágil con entregables incrementales, priorización de módulos
críticos y cronograma con holgura. \\
R6 & Transferir & Proveedor cloud confiable con SLA, backups automáticos y plan de
recuperación; respaldo local opcional. \\
R7 & Mitigar & Autenticación por roles, cifrado, VPN WireGuard para acceso a MikroTik
y buenas prácticas de seguridad. \\
\bottomrule
\end{tabularx}
\end{center}

% ======================================================================
\section{Gestión Ágil (Scrum)}
\label{sec:scrum}

Para la construcción del software se adopta el marco \textbf{Scrum}: el trabajo se
organiza en un Product Backlog priorizado, se planifica y ejecuta en Sprints de
duración fija, y cada Sprint entrega un Incremento potencialmente utilizable. Este
enfoque adaptativo es reconocido por el PMBOK 7.ª edición para proyectos con requisitos
que evolucionan.

\subsection{Roles del equipo Scrum}

\begin{center}
\captionof{table}{Roles del equipo Scrum.}
\label{tab:roles-scrum}
\begin{tabularx}{\textwidth}{L{2.6cm} L{3.6cm} X}
\toprule
\textbf{Rol Scrum} & \textbf{Integrante} & \textbf{Responsabilidad} \\
\midrule
Product Owner & Cristian Cabana Sulca & Define y prioriza el backlog; representa al
negocio (conoce a WISP); valida los incrementos. \\
Scrum Master & Jherson J. P. Fernández S. & Facilita las ceremonias, remueve
impedimentos y protege al equipo. \\
Developer & Alessandro Pastor Mamani M. & Desarrollo backend y aplicación móvil. \\
Developer & Andrés Lino Montes Mamani & Desarrollo backend y frontend web. \\
\bottomrule
\end{tabularx}
\end{center}

\subsection{Product Backlog priorizado}

El Product Backlog recoge las historias de usuario priorizadas por valor. Se muestran
las principales, con su prioridad y estimación en puntos de historia
(\textit{Story Points}):

\begin{center}
\captionof{table}{Product Backlog priorizado (principales historias de usuario).}
\label{tab:backlog}
\begin{tabularx}{\textwidth}{L{1.6cm} X C{1.8cm} C{1cm}}
\toprule
\textbf{ID} & \textbf{Historia de usuario} & \textbf{Prioridad} & \textbf{SP} \\
\midrule
HU-01 & Como vendedor, quiero registrar prospectos desde la app móvil para no usar
papel ni WhatsApp. & Alta & 8 \\
HU-02 & Como administrador, quiero validar prospectos y generar órdenes de
instalación. & Alta & 5 \\
HU-03 & Como técnico, quiero registrar la instalación (fotos NAP, borne, equipos)
desde la app. & Alta & 8 \\
HU-04 & Como administrador, quiero crear/cortar/reactivar clientes PPPoE en el
MikroTik desde el sistema. & Alta & 13 \\
HU-05 & Como administrador, quiero gestionar clientes, planes y zonas desde el panel
web. & Media & 8 \\
HU-06 & Como visitante, quiero ver la cobertura de la empresa en el sitio web. &
Media & 5 \\
HU-07 & Como cliente, quiero recibir notificaciones automáticas por WhatsApp. &
Baja & 5 \\
\bottomrule
\end{tabularx}
\end{center}

\subsection{Sprint Planning (planificación de sprints)}

El desarrollo se organiza en sprints de 2 semanas. La documentación de GTI (semestre~1)
se distribuye en los primeros sprints, y la construcción del software (semestre~2) en
sprints posteriores. Planificación de referencia:

\begin{center}
\captionof{table}{Planificación de sprints de referencia.}
\label{tab:sprints}
\begin{tabularx}{\textwidth}{L{2.2cm} L{2.8cm} X}
\toprule
\textbf{Sprint} & \textbf{Fechas} & \textbf{Foco / Incremento} \\
\midrule
Sprint 1 & 16/03 -- 29/03 & Inducción, diagnóstico organizacional. \\
Sprint 2 & 30/03 -- 12/04 & Business Case, inicio del plan de gestión. \\
Sprint 3 & 13/04 -- 26/04 & Plan de gestión, procesos AS-IS. \\
Sprint 4 & 27/04 -- 10/05 & Procesos TO-BE, inicio SRS. \\
Sprint 5 & 11/05 -- 24/05 & SRS y prototipos (Brief aprobado). \\
Sprint 6 & 25/05 -- 07/06 & Arquitectura TIC, UML. \\
Sprint 7 & 08/06 -- 21/06 & Diseño de base de datos, consolidación. \\
Sprints 8+ & Semestre 2 & Construcción: backend + MikroTik, frontend, móvil, WhatsApp,
pruebas, despliegue. \\
\bottomrule
\end{tabularx}
\end{center}

\subsection{Incrementos y Definición de Terminado (DoD)}

Cada sprint produce un incremento que cumple la Definición de Terminado
(\textit{Definition of Done}) del proyecto:

\begin{itemize}[leftmargin=1.4em]
  \item El entregable/funcionalidad está completo según su criterio de aceptación.
  \item Ha sido revisado por el Product Owner (documentos) o probado (software).
  \item Está versionado en el repositorio Git (para el software).
  \item Cumple los estándares de calidad definidos (IEEE/ISO según corresponda).
\end{itemize}

\subsection{Mecanismos de control y seguimiento}

El control del proyecto se realiza mediante las ceremonias Scrum y el seguimiento del
avance:

\begin{itemize}[leftmargin=1.4em]
  \item \textbf{Daily / reuniones de sincronización:} coordinación frecuente del
  equipo sobre avance e impedimentos.
  \item \textbf{Sprint Review:} al cierre de cada sprint se presenta el incremento al
  Product Owner y, en hitos, a los revisores.
  \item \textbf{Sprint Retrospective:} mejora continua del proceso del equipo.
  \item \textbf{Revisiones de avance (hitos H2, H4):} control formal ante la escuela
  en las semanas 5-6 y 11-12.
  \item \textbf{Tablero Kanban / gestión visual:} seguimiento del estado de las tareas
  (Por hacer / En progreso / Hecho).
\end{itemize}

% ======================================================================
\begin{upeubox}{Conclusión del Plan de Gestión}
El presente Plan de Gestión establece la hoja de ruta completa del proyecto bajo un
enfoque híbrido PMBOK + Scrum. Define formalmente el proyecto (acta de constitución),
delimita su trabajo (EDT y diccionario), programa sus actividades (cronograma, Gantt e
hitos), establece su línea base de costos, gestiona sus riesgos con un plan de
respuesta, y organiza la construcción del software bajo Scrum (backlog, sprints e
incrementos).

\medskip
\textbf{Trazabilidad y disciplina metodológica.} Cada componente se alinea a su marco
de referencia: el acta al \textit{Project Charter} del PMI, la EDT a la regla del
100\,\%, el cronograma a la gestión del tiempo con hitos, el presupuesto a la línea
base de costos, los riesgos al análisis cualitativo $P\times I$ con estrategias de
respuesta, y la gestión ágil a Scrum. Con ello, el proyecto cuenta con la disciplina
metodológica y la trazabilidad de resultados exigidas, y queda listo para su ejecución.
\end{upeubox}
